% Suggested changes to chapters/06_ApplicationWindFarmControl_04.tex,
% section "WFSim Program Structure and Code Modifications": text matching
% the updated Figure fig:WFSimFlow(b) (FlowchartInitCtrlPlant_2), which now
% shows the current code (WFSim_demo, default closed loop: WRC-AIC with the
% rotor effective wind speeds estimated from the power).
% Each block: the original line(s) as "% OLD (line n):", followed by the
% new text. Unchanged lines are not repeated.

% =====================================================================
% OLD (line 24, last two sentences):
% The diagrams are generated using the MATLAB profiler during a closed-loop run with the nonlinear MPC controller integrated into WFSim and the Koopman MPC published in \cite{sharan2022Koopman}. The functions are reordered based on their dependencies to visualize the sequence of function calls.
% NEW:
The diagrams are generated using the MATLAB profiler during a closed-loop run with the nonlinear MPC controller integrated into WFSim and with the WRC-AIC Koopman MPC of this thesis, see Section~\ref{sec:WindFarmLinearMPCforWRC-AIC}. The functions are reordered based on their dependencies to visualize the sequence of function calls.

% =====================================================================
% OLD (line 56):
%     \item In the modified program flow, two new functions are added to the initialization phase: \textit{MPCinitK.m} and \textit{koopmanstateextensionWFSim.m}. These functions load the state-space model containing the finite-dimensional representation of the Koopman operator, which is used for the subsequent MPC execution.
% NEW:
    \item In the modified program flow, two new functions are added to the initialization phase: \textit{MPCinitv.m} and \textit{koopmanstateextensionWFSim.m}. The function \textit{MPCinitv.m} sets all constant controller parameters and loads the state-space models containing the finite-dimensional representation of the Koopman operator, which are used for the subsequent MPC execution. The function \textit{koopmanstateextensionWFSim.m} evaluates the Koopman lifting functions to compute the initial lifted state.

% =====================================================================
% OLD (line 80):
% In the modified structure (Figure~\ref{fig:WFSimFlow}b), this workflow is replaced by the \textit{qLPVMPCcontroller\_dU\_\allowbreak quadprog.m} function. This new controller utilizes matrices initialized in \textit{koopmanstateextensionWFSim.m} and initializes wind-dependent input matrices from \textit{MPCinit\_ctrl.m}. The functions \textit{wfmodel.m} and \textit{matsys.m} are modified to accommodate these changes and to construct the parameter-dependent Toeplitz-like matrices and boundary conditions for the QP formulation. These are passed to \textit{quadprog\_sol.m}, which calls the MATLAB solver \textit{quadprog.m}. While the original implementation performed two iterations to refine wind speed predictions, no clear benefit regarding tracking error minimization was observed in the studies forming the basis of this thesis when running the optimization twice. The loop with two iterations is retained in the code. However, as the predicted effective wind speeds do not depend on the decision variables, the second iteration yields the same solution as the first one, so that effectively the farm control signal is calculated only once at each sample.
% NEW:
In the modified structure (Figure~\ref{fig:WFSimFlow}b), this workflow is replaced by the \textit{qLPVMPCcontroller\_dU\_\allowbreak quadprog.m} function. Before the controller is called, \textit{urFromPower.m} estimates the rotor effective wind speeds of the turbines from the measured power, the thrust coefficient $C^{\prime}_T$ and the yaw angle $\gamma$ by inverting the power model of \textit{Actuator.m}, so that the controller does not use wind speeds measured in WFSim. The controller utilizes the Koopman models and constant matrices initialized in \textit{MPCinitv.m} and updates the wind-dependent input matrices in \textit{MPCinit\_ctrl.m}. If the power demand exceeds the power available with greedy control, \textit{findMaxYaw.m} determines the yaw angle of the upstream turbine that maximizes the farm power according to an analytical wake model, which serves as the yaw reference of the WRC. The function \textit{wfmodelNL.m}, derived from \textit{wfmodel.m}, predicts the rotor effective wind speeds over the horizon with the Koopman model and forms the wind farm system and input matrices. Based on these, the modified function \textit{Matsys.m} constructs the parameter-dependent Toeplitz-like prediction matrices, and \textit{MatsysP.m} builds the corresponding prediction matrices of the Koopman power model with the inputs $C^{\prime}_{T,1}$, $C^{\prime}_{T,2}$ and $\gamma$. These are passed to \textit{quadprog\_sol.m}, which computes the Hessian and gradient of the QP in \textit{H\_G\_matrix\_dU.m} and, for the Koopman power model, in \textit{H\_G\_Kmatrix\_dU.m}, and calls the MATLAB solver \textit{quadprog.m}. While the original implementation performed two iterations to refine wind speed predictions, no clear benefit regarding tracking error minimization was observed in the studies forming the basis of this thesis when running the optimization twice. The loop with two iterations is retained in the code. However, as the predicted effective wind speeds do not depend on the decision variables, the second iteration yields the same solution as the first one, so that effectively the farm control signal is calculated only once at each sample.
